\section{Related Work}
\label{sec:related}

\paragraph{4D Scene Reconstruction}
Early 4D reconstruction works mainly focus on optimization-based approaches~\cite{newcombe2015dynamicfusion,pumarola2021d,li2021neural,du2021neural,tretschk2021non,park2021nerfies,fridovich2023k,cao2023hexplane,li2023dynibar,Wu_2024_CVPR,yang2024deformable,yang2024real,som2024,wang2025freetimegs},
which iteratively fits a 4D representation to monocular or multi-view videos.
With the development of neural radiance fields (NeRFs)~\cite{mildenhall2020nerf},
many time-dependent NeRFs~\cite{park2021nerfies,du2021neural,li2021neural,pumarola2021d,fridovich2023k,cao2023hexplane,li2023dynibar}
fit deformable 3D representations to dynamic scenes.
However, these approaches suffer from expensive volumetric rendering,
making them less practical for real-world applications.
3D Gaussian Splatting (3D-GS)~\cite{kerbl20233d} avoids expensive sampling using a rasterization-based rendering pipeline.
Several works~\cite{Wu_2024_CVPR,yang2024deformable,yang2024real,som2024,zhu2024motiongs,lu2024dn,wang2025freetimegs}
extend it to dynamic scene reconstruction,
which significantly reduces the rendering time.
Besides,
some methods~\cite{Zhang2022StructureAM,li2025_megasam,lu2025trackingworld} achieve accurate and robust 4D reconstruction by leveraging depth priors~\cite{hu2024-DepthCrafter,ke2024repurposing,piccinelli2024unidepth,yang2024depth,lihe24DAV2}.
However, they still require per-scene optimization.

Some recent works~\cite{zhang24monst3r,jin2024stereo4d,wang2025continuous,sucar2025dynamic,st4rtrack2025,jiang2025geo4d,xu20254dgt,AETHER,chen2025back,liang2024feed} have explored feed-forward 4D scene reconstruction from monocular videos.
Among them,
MonST3R~\cite{zhang24monst3r} adapts the notable \emph{static} 3D reconstructor DUSt3R~\cite{wang2024dust3r} to dynamic scenes.
The follow-ups~\cite{jin2024stereo4d,wang2025continuous,sucar2025dynamic,st4rtrack2025,AETHER,lu2025align3r}
took a similar path, explicitly predicting the point correspondences.
However, due to DUSt3R's limitations, they process pairs of frames at a time.
To handle long monocular videos,
$\pi^3$~\cite{wang2025pi} builds a permutation-equivariant architecture on top of VGGT~\cite{wang2025vggt} for static and dynamic 3D reconstruction.
Geo4D~\cite{jiang2025geo4d} leverages video generators~\cite{xing2024dynamicrafter} to directly infer 4D point maps from monocular videos.
4DGT~\cite{xu20254dgt} and BTimer~\cite{liang2024feed} utilize transformer-based architectures to predict dynamic 3D Gaussian representations~\cite{kerbl20233d}.
However, they do not explicitly model dense point correspondences over time.

\paragraph{Scene Flow Estimation}
Early works define point correspondences as optical flow estimation in pixel space~\cite{horn1981determining,brox2004high,bruhn2005lucas,ilg2017flownet,teed2020raft}.
One popular pipeline is to directly estimate dense pixel-wise correspondences in a coarse-to-fine manner~\cite{dosovitskiy2015flownet,ilg2017flownet,sun2018pwc,teed2020raft}.
However, such a coarse-to-fine strategy may fail in the presence of large motions or occlusions~\cite{revaud2015epicflow}.
More recently,
GMFlow~\cite{Xu_2022_CVPR,xu2023unifying} proposes to reformulate optical flow estimation as a global matching problem rather than local regression.
The follow-ups~\cite{huang2022flowformer,shi2023flowformer++} also use transformer-based neural networks to model the global correlations.
% As for the long-range point tracking~\cite{rubinstein2012towards,sand2008particle},
% recent works~\cite{doersch2022tap,karaev2024cotracker,karaev2025cotracker3} explore to learn robust point trackers from large-scale video datasets.
However, these methods still deal with 2D point correspondences in image space.
In contrast,
we address 3D scene flow estimation in world space.
Some researchers have also explored estimating 3D \emph{scene flow} directly from image pairs,
including RAFT-3D~\cite{teed2021raft}, SpatialTracker~\cite{SpatialTracker}, SceneTracker~\cite{wang2025scenetracker}, and TAPVid-3D~\cite{koppula2024tapvid}.
More closely related to our work,
several recent works~\cite{sucar2025dynamic,st4rtrack2025,jin2024stereo4d} explore reconstructing dynamic 3D geometry,
along with 3D scene flow estimation in world space.
However, they process only two images at a time and require post-processing to refine the results.

\paragraph{Geometric Diffusion Model}
Like our approach,
many recent works have leveraged pre-trained off-the-shelf diffusion models~\cite{rombach2022high,ho2022video,wang2023modelscope,singer2023makeavideo,blattmann2023align,ge2023preserve,wang2023videocomposer,guo2024animatediff,blattmann2023stable,zhang2024show,huang2024blue,xing2024dynamicrafter,kong2024hunyuanvideo,liang2025diffusion,garcia2025fine}
to tackle 3D tasks~\cite{gao2024cat3d,sargent2024zeronvs,chen2024mvsplat360,yu2024viewcrafter,szymanowicz2025bolt3d,li2025dso,wu2025amodal3r,ke2024repurposing,ye2024stablenormal,xu2024matters,song2025depthmaster,liu2023zero,shi2024mvdream,zheng2024free3d,voleti2025sv3d}, thanks to rich priors learned from large-scale image or video datasets.
When it comes to 4D reconstruction, a straightforward solution is to generate multi-view videos,
and then fit a 4D representation via per-scene optimization~\cite{xie2024sv4d,zeng2024stag4d,wu2025cat4d}.
Inspired by score distillation sampling (SDS)~\cite{pooledreamfusion},
another line of works~\cite{jiang2024consistentd,zhang20244diffusion,dreamscene4d,ren2023dreamgaussian4d,li2024dreammesh4d}
directly distill 4D priors from pre-trained video generators.
However, these approaches still rely on iterative per-scene optimization,
which is expensive when dealing with in-the-wild videos.
The most related work is Geo4D~\cite{jiang2025geo4d} that fine-tunes a pre-trained video diffusion model to directly infer dynamic 3D point maps, depths, and camera poses from monocular videos.
Our method differs from Geo4D in two aspects:
(1) we \emph{simultaneously} reconstruct \emph{dynamic 3D geometry} and estimate \emph{dense point correspondences} in a unified 4D VAE framework; and
(2) we show that it is \emph{not} necessary to align the data and latent spaces when fine-tuning a diffusion model.